
\clearpage
\chapter{סיכום ועתיד האבטחה בעידן הסוכנים}

\section{תרומות}

מאמר זה הציג את מודל האיום \textenglish{ClawHavoc} ומסגרת ההגנה \textenglish{SkillSieve} עבור שרשרת האספקה בצינורות סוכן-\textenglish{LLM}. ה-\textenglish{SkillSieve} מדגים \textenglish{F1 = 0.92} על מערך בדיקה מגוון, עם עלות ביצוע נסבלת לייצור.

\section{המלצות מעשיות}

\begin{itemize}
\item כל מסגרת סוכן חייבת לממש שכבת אימות כישורים לפני טעינה (\textenglish{pre-load validation}).
\item ריפוזיטורי כישורים ציבוריים זקוקים למנגנון דיווח שקוף על פגיעויות.
\item ביקורת \textenglish{audit log} של כל פעולות הכישורים הכרחית לזיהוי ומניעת פרצות עתידיות.
\end{itemize}

\section{כיוונים עתידיים}

מחקר עתידי יתמקד בזיהוי תקיפות מרובות שלבים (\textenglish{multi-step attacks}) בהן כישורים לגיטימיים לכאורה משתלשלים ליצור שרשרת זדונית. כמו כן, שילוב עם \textenglish{formal verification} לכישורים קריטיים עשוי לאפשר ערבויות אבטחה פורמליות \cite{weidinger2021ethical}.
